%%%BLOCK id=081-01 ch=99 sec=页边零星·化学 kind=misc alt=none cont=none y=0.06-0.10
% 页面上缘空白处(页眉“No.”旁)有零星字迹，疑为相邻页内容(化学·气体密度)，照录
$\rho=1.429\unit{g\cdot L^{-1}}$
%%%END
%%%BLOCK id=083-01 ch=99 sec=页边零星·化学 kind=misc alt=none cont=none y=0.03-0.08
% 页面上缘右侧页眉“No.”旁有零星字迹(与p081上缘相同，疑为化学·气体密度)，照录
$\rho=1.429\unit{g\cdot L^{-1}}$
%%%END
%%%BLOCK id=085-01 ch=99 sec=页边零星·化学 kind=misc alt=none cont=none y=0.02-0.07
% 页面上缘右侧页眉“No.”旁有零星字迹(与p081、p083上缘相同，疑为化学·气体密度)，照录
$\rho=1.429\unit{g\cdot L^{-1}}$
%%%END
%%%BLOCK id=127-02 ch=99 sec=化学·气体密度 kind=note alt=none cont=none y=0.05-0.09
% 页面顶部右侧，压在页眉线上，右端被页边截断
\unclear{气}体 $\rho=1.429\unit{g\cdot L^{-1}}$
%%%END
%%%BLOCK id=129-02 ch=99 sec=化学·气体密度 kind=note alt=none cont=none y=0.05-0.09
% 页面顶部右侧，压在页眉线上，右端被页边截断；与第127页同一位置出现同一行字
\unclear{气}体 $\rho=1.429\unit{g\cdot L^{-1}}$
%%%END
%%%BLOCK id=131-02 ch=99 sec=化学·气体密度 kind=note alt=none cont=none y=0.03-0.06
% 页面右上角，压在页眉线上，绝大部分被页边截断，仅可见末尾（与第127、129页同一位置的同一行字）
\unclear{$\ldots$}\,$\unit{g\cdot L^{-1}}$
%%%END
%%%BLOCK id=131-08 ch=99 sec=涂鸦·无关数字 kind=misc alt=none cont=none y=0.88-0.95
% 圆圈内的数字（每个下方有下划线与小点）：100、50、67、150，其后另有一个被箭头划去的 150；其下有一排涂划线（像连写的“6”）
\unclear{100\ \ 50\ \ 67\ \ 150}
% 另：页面左侧(“W”形线)、右下角(网格状)有大片涂划，未转录
%%%END
%%%BLOCK id=131-09 ch=99 sec=数学·函数 kind=misc alt=none cont=none y=0.94-1.0
% 页面左下角，旁有波浪线涂划
\[ \unclear{f(m)}=\dfrac{m^{2}+2m+\unclear{1}}{2m+3} \]
%%%END
%%%BLOCK id=133-06 ch=99 sec=化学·零星 kind=note alt=none cont=none y=0.03-0.07
% 页眉右上角“No.”处有一行淡灰色字迹（疑为相邻页内容；p135 顶部同一位置也有同样一行），与物理无关
\unclear{混合气体} $\rho=1.429\,\unit{g\cdot L^{-1}}$
%%%END
%%%BLOCK id=135-06 ch=99 sec=化学·零星 kind=note alt=none cont=none y=0.02-0.07
% 页眉右上角“Date.”处有一行淡灰色字迹（疑为相邻页内容；p133 顶部同一位置也有同样一行），与物理无关
\unclear{混合气体} $\rho=1.429\,\unit{g\cdot L^{-1}}$\quad（右端另有一处字迹 \unclear{$-6.72$}）
%%%END
%%%BLOCK id=169-08 ch=99 sec=涂鸦·试笔波浪线 kind=misc alt=none cont=none y=0.66-0.97
% 页面右侧大片试笔涂画：多条波浪线，其间写有两处英文 math；另在 y≈0.68 处有几笔零散折线
涂鸦：多条波浪线，其间写有 math、math
%%%END
%%%BLOCK id=181-01 ch=99 sec=化学·混合气体计算 kind=misc alt=none cont=none y=0.00-0.11
% 页面右上角叠压着另一页（带 Date. NO. 页眉），这几行是化学笔记，字迹倾斜较乱
\unclear{0}.4mol\quad $\mathrm{CO}$\quad $\mathrm{CO_2}$

混合气 充分燃烧得 \unclear{O$_2$} 和 $\mathrm{CO_2}$ 混合气体 $\rho=1.429\,\mathrm{g/L}$

$\mathrm{CO}$为 $0.4\times\dfrac{3}{4}\times22.4=6.72\,\mathrm{L}$

\unclear{…$h\rho$\,g/mL}的密度 % 右侧叠压页上的字，位置在“不用测k,m”一行的右上方，字迹很乱
%%%END
%%%BLOCK id=196-02 ch=99 sec=化学·吸附态相对能量(非物理内容) kind=note alt=none cont=none y=0.31-0.34
% 与物理无关的化学笔记（含吸附态记号 $*$）；“相对能量”四字为写在 CH$_3$OH$^*$ 上方的小字
% 此行左下方(y约0.355-0.37,x约0.02-0.06)另有一处小记号(一道竖线加两三道斜线，含义不明，疑似涂鸦)
CH$_3$OH$^{*}$\textsuperscript{相对能量}\ 不是甲醇的能量\qquad CO$^{*}$\ 能量不是CO能量
%%%END
%%%BLOCK id=199-03 ch=99 sec=生物·标志重捕法与稀释涂布平板计数 kind=note alt=none cont=none y=0.27-0.42
% 与物理无关的生物笔记；后三行的“$\rho$小”被同一个椭圆圈起，首行“$\rho$大”下有一道横线
\begin{itemize}
\item 标志重捕法\ 标志物脱落\quad $\rho$大
\item 计数酵母菌时未稀释\quad $\rho$小
\item 在稀疏区取样\quad $\rho$小
\item 用稀释涂布平板计数大肠杆菌\quad $\rho$小
\end{itemize}
%%%END
%%%BLOCK id=199-04 ch=99 sec=化学·有机反应(内成环加成) kind=note alt=none cont=none y=0.45-0.50
% 与物理无关的化学笔记；红笔下划线共三处（“内成环”、“也”、“加成反应”，文字本身为黑色）
分子\red{\uline{内成环}}\red{\uline{也}}叫\red{\uline{加成反应}}，（啥也没加）。
%%%END
%%%BLOCK id=199-05 ch=99 sec=化学·氢化物与硅硫硼化合物 kind=note alt=none cont=none y=0.51-0.69
% 与物理无关的化学笔记。红笔：首句“氢化”下、方框内“最简单氢化物”下、两个圆圈内文字下各有一道红色下划线；方框内压着一个五角星
\red{\uline{氢化}}物注意\ \fbox{最简单氢化物 ★}\ \ \fbox{CH$_4$}\ 与\ \fbox{\unclear{沥青}}\ 都是氢化物

\notefig[0.6]{figures/p199_d.png}
% 结构式转录：(CH$_3$)$_3$ *Si$-$S$-$B，B 上分两支：上支 $-$S$-$Si(CH$_3$)$_3$；下支 S \struck{CH$_3$} 另起 Si$-$(CH$_3$)$_3$
%%%END
%%%BLOCK id=200-01 ch=99 sec=生物·微生物高通量测序 kind=misc alt=none cont=none y=0.00-0.09
% 页首粘贴的紫色纸条（生物题），上缘被截断；灰色印刷字与黑色手写字叠在一起
\textbf{（页首粘贴的紫色纸条，生物题；印刷字与手写字叠在一起，上缘被截断）}

手写（纸条最上一行）：测序法\unclear{与对核酸}序列，可快速区分新微生物种类

印刷：“活性污泥”中的微生物种类进行研究时，常采取高通量 DNA 测序手段，\unclear{……}分离培养手段\unclear{……}你推测原因可能是

手写：④ 在短时间内进行大量DNA片断测序，① 未知的微生物种类太多，不知道该如何\unclear{培养}；② 某些微生物以\unclear{较多}存在，分离培养上操作\unclear{困难}
%%%END
%%%BLOCK id=207-05 ch=99 sec=生物·果蝇性染色体检测 kind=misc alt=none cont=none y=0.76-0.99
\textbf{检测雄果蝇性染色体组成：镜检}

（下方粘贴的印刷纸条，生物题答案）(3)取雄果蝇造血干细胞（或其他能观察细胞分裂的细胞）做成装片，在显微镜下观察细胞有丝分裂中期的（Y/性）染色体数目\ （2分，合理即可）若有丝分裂中期染色体数目为8条（Y染色体数目为1条/性染色体数目为2条），则可育雄果蝇的性染色体组成为XY；若有丝分裂中期染色体数目为9条（Y染色体数目为2条/性染色体数目为3条），则可育雄果蝇的性染色体组成为XYY\ （4分，合理即可）
%%%END
%%%BLOCK id=207-06 ch=99 sec=生物·微生物培养避光 kind=note alt=none cont=none y=0.85-0.91
微生物培养时避光\unclear{培养}目的：排除光降解对实验结果的影响
%%%END
%%%BLOCK id=209-02 ch=99 sec=化学·化学平衡与平衡常数 kind=note alt=none cont=none y=0.62-0.93
左=右的反应 有连锁反应时 $P\uparrow$ 反应平衡会移动

$K_{\text{总}}=\dfrac{K_1}{K_2}$\qquad $K_{\text{总}}=K_1\cdot K_2$

\hspace{4.2cm}$K_1=\dfrac{K_{\text{总}}}{K_2}$
%FIG p209_b 0.46 0.71 0.66 0.93
\notefig[0.25]{figures/p209_b.png}
%%%END
%%%BLOCK id=297-02 ch=99 sec=纸条残片 kind=misc alt=none cont=none y=0.57-0.85
页面左缘一条灰色纸条上露出的印刷字，自上而下：“.”、向、\unclear{愿}、\unclear{身}、夹、在、A.、C.（疑为第295页所贴印刷题各行行首字“向／感／射／夹／在／A.／C.”的残边，无其它内容）
%%%END
